\section{Computing policies and hardest records on polytopes}\label{sec:general-algorithms}
Let $F=\{y\in\mathbb R^n:Ay\le b\}$ be nonempty and compact, and let
$T(y)=c^\top y+d$. Equalities can be represented by paired inequalities.
The constant $d$ does not affect width. The first formulation allows any
declared finite family $\mathcal Q$ of legal batches; $L_Q$ denotes the matrix
of their exact linear measurements.

\subsection{An enumerative hardest-record formulation}
Introduce a common record $y$, endpoint witnesses $u_Q,v_Q$ for each
$Q\in\mathcal Q$, and a scalar $\eta$. Consider
\begin{equation}\label{eq:general-hardest-lp}
\begin{aligned}
\max\quad &\eta\\
\text{subject to}\quad &Ay\le b,\quad Au_Q\le b,\quad Av_Q\le b
                         &&(Q\in\mathcal Q),\\
&L_Qu_Q=L_Qy,\quad L_Qv_Q=L_Qy &&(Q\in\mathcal Q),\\
&\eta\le c^\top(u_Q-v_Q) &&(Q\in\mathcal Q).
\end{aligned}
\end{equation}
For fixed $y$, the independent witnesses for batch $Q$ can attain its sharp
width $W_Q(y)$. Thus the largest feasible $\eta$ at that record is
$\min_Q W_Q(y)$, and the optimum of \eqref{eq:general-hardest-lp} is $\rho$.
Its common record is a certificate that every legal batch leaves width at
least $\rho$ there.

The nonnegative dual multipliers $\pi_Q$ of the last constraints form an
optimal randomized batch policy. Indeed, the free variable $\eta$ forces
$\sum_Q\pi_Q=1$. Dualizing these constraints first yields
\[
\min_{\pi\in\Delta(\mathcal Q)}
 \max_{y,u_Q,v_Q}\sum_Q\pi_Qc^\top(u_Q-v_Q)
 =\min_{\pi\in\Delta(\mathcal Q)}\max_{y\in F}
       \sum_Q\pi_Q W_Q(y)=\rho,
\]
where the endpoint feasibility constraints remain in the inner maximization.
Ordinary linear-programming strong duality supplies the equality and the
policy. In the implementation, a separate
LP fixes the extracted $\pi$ and maximizes its expected width, checking the
upper bound independently of the constraints involving $\eta$.

With $M=|\mathcal Q|$ batches, the formulation has $n(1+2M)+1$ variables.
It is polynomial in the explicitly enumerated batch family, which may be
exponential in the number of available measurements. It is a small-instance
benchmark and certificate generator; the aggregate-specific compact
formulation avoids this enumeration in its structured setting.

\subsection{An exact deterministic two-query algorithm}\label{sec:two-query}
Assume every one of the $n\ge2$ coordinates is queryable at unit cost and the
hard budget is two. An unused query can be added without increasing width,
so considering two distinct coordinates includes policies that stop earlier.
For a first query $i$ and feasible answer $t$, define
\begin{equation}\label{eq:conditional-two-query}
g_{ij}(t)=\max\{c^\top(u-v):Au\le b,\ Av\le b,\
                u_i=v_i=t,\ u_j=v_j\},\qquad j\ne i.
\end{equation}
This maximizes the terminal width over all feasible records having first
answer $t$, after second query $j$. The best deterministic policy beginning
with $i$ therefore has worst width
$\rho_i=\max_t\min_{j\ne i}g_{ij}(t)$.

\begin{proposition}\label{prop:two-query-lp}
For each $i$, $\rho_i$ is the value of the following single linear program,
with one endpoint pair $(u^j,v^j)$ for each $j\ne i$:
\begin{equation}\label{eq:two-query-lp}
\begin{aligned}
\max\quad &\eta\\
\text{subject to}\quad &Au^j\le b,\quad Av^j\le b &&(j\ne i),\\
&(u^j)_i=(v^j)_i=t,\quad (u^j)_j=(v^j)_j &&(j\ne i),\\
&\eta\le c^\top(u^j-v^j) &&(j\ne i).
\end{aligned}
\end{equation}
The optimal deterministic adaptive two-query worst-width value is
$\min_i\rho_i$.
\end{proposition}
\begin{proof}
At fixed feasible $t$, each endpoint pair has only its own constraints and
can attain $g_{ij}(t)$. All inequalities on $\eta$ are feasible precisely up
to their smallest attainable upper bound, $\min_{j\ne i}g_{ij}(t)$.
Maximizing over the shared $t$ proves the first claim.
For any deterministic policy beginning with $i$, its chosen second coordinate
at answer $t$ has conditional worst width at least $\min_jg_{ij}(t)$.
Conversely, choosing a minimizing $j$ after observing $t$ attains that value.
Finite tie breaking gives a measurable selection because these polyhedral
value functions are piecewise linear on their feasible interval. Finally,
minimize over the first coordinate.
\end{proof}

For rational input, $n$ LPs of polynomial description size choose the first
coordinate; after its answer, $n-1$ conditional LPs implement the second
choice. This is a polynomial-time deterministic two-query procedure in the
standard bit-complexity model, including the encoding length of rational
query answers. For arbitrary exact answers, the policy is still defined by
the conditional LPs; the polynomial bit-complexity statement applies to
rationally encoded inputs and answers. It does not solve arbitrary-budget
randomized precision-risk optimization or minimize expected requests until
certification. No grid of possible first answers is used.
